% This is the list of proof rules

\begin{proof-rule}{assume}{veriT}
\begin{plContainer}
\begin{plList}
\proofsep& i.& \phi &(\currule)\\
\end{plList}
\end{plContainer}
where $\varphi$ corresponds to a formula asserted in the input problem up
to the orientation of equalities.
\paragraph{Remark.} 
This rule can not be used by the
$\grT{(step }\dots\grT{)}$ command. Instead it corresponds to the dedicated
$\grT{(assume }\dots\grT{)}$ command.

\end{proof-rule}

\begin{proof-rule}{true}{veriT}
\begin{plContainer}
\begin{plList}
\proofsep& i.& \top  &(\currule)\\
\end{plList}
\end{plContainer}
\end{proof-rule}

\begin{proof-rule}{false}{veriT}
\begin{plContainer}
\begin{plList}
\proofsep& i.& \neg \bot &(\currule)\\
\end{plList}
\end{plContainer}
\end{proof-rule}

\begin{proof-rule}{not_not}{veriT}
\begin{plContainer}
\begin{plList}
\proofsep& i.& \neg (\neg \neg \varphi) , \varphi &(\currule)\\
\end{plList}
\end{plContainer}

\paragraph{Remark.} This rule is useful to remove double negations from a
clause by resolving a clause with the double negation on \(\varphi\).

\end{proof-rule}

\begin{proof-rule}{th_resolution}{veriT}
This rule is the resolution of two or more clauses.

\begin{plContainer}
\begin{plList}
\proofsep& i_1.& \varphi^{1}_{1} , \dots , \varphi^{1}_{k^1} &(\dots)\\
\plLine\\
\proofsep& i_n.& \varphi^{n}_{1} , \dots , \varphi^{n}_{k^n} &(\dots)\\
\proofsep& j.&
\varphi^{r_1}_{s_1} , \dots , \varphi^{r_m}_{s_m} 
&(\currule; i_1, \dots, i_n)\\
\end{plList}
\end{plContainer}
where $\varphi^{r_1}_{s_1} , \dots , \varphi^{r_m}_{s_m}$ are from $\varphi^{i}_{j}$ and
are the result of a chain of predicate resolution steps on the clauses of
$i_1$ to $i_n$. It is possible that $m = 0$, i.e. that
the result is the empty clause.

This rule is only used when the resolution step is not emitted by the SAT solver.
See the equivalent \proofRule{resolution} rule for the rule emitted by the
SAT solver.

\paragraph{Remark.} While checking this rule is NP-complete, the \currule-steps
produced by veriT are simple. Experience with reconstructing the step in
Isabelle/HOL shows that checking can done by naive decision procedures. The
vast majority of \currule-steps are binary resolution steps.
\end{proof-rule}

\begin{proof-rule}{resolution}{veriT}
This rule is equivalent to the \proofRule{the\_resolution} rule, but it is
emitted by the SAT solver instead of theory reasoners. The differentiation
serves only informational purpose.

\end{proof-rule}

\begin{proof-rule}{tautology}{veriT}
  \begin{plContainer}
    \begin{plList}
      \proofsep& i.&
      \varphi_1 , \dots , \varphi_i , \dots , \varphi_j ,\dots , \varphi_n
      &(\dots)\\
      \proofsep& j.& \top &(\currule; i)\\
    \end{plList}
  \end{plContainer}
  and $\varphi_i = \bar\varphi_j$.
\end{proof-rule}

\begin{proof-rule}{contraction}{veriT}
  \begin{plContainer}
    \begin{plList}
      \proofsep& i.& \varphi_1 , \dots , \varphi_n &(\dots)\\
      \plLine\\
      \proofsep& j.& \varphi_{k_1} , \dots , \varphi_{k_m}
      &(\currule; i)\\
    \end{plList}
  \end{plContainer}
  where $m \leq n$ and  $k_1 \dots k_m$ is a monotonic map to $1 \dots n$
  such that $\varphi_{k_1} \dots \varphi_{k_m}$ are pairwise distinct and
  $\{\varphi_1, \dots, \varphi_n\} = \{\varphi_{k_1} \dots \varphi_{k_m}\}$.
  Hence, this rule removes duplicated literals.
\end{proof-rule}

\begin{proof-rule}{subproof}{veriT}
The \currule{} rule completes a subproof and discharges local
assumptions. Every subproof starts with some \proofRule{input} steps. The
last step of the subproof is the conclusion.

\begin{plContainer}
\begin{plSubList}
\proofsep& i_1.& \psi_1 &(\proofRule{input})\\
\plLine\\
\proofsep& i_n.& \psi_n &(\proofRule{input})\\
\plLine\\
\proofsep& j.& \varphi &(\dots)\\
\end{plSubList}
\begin{plList}
\proofsep& k.& \neg\psi_1 , \dots , \neg \psi_n ,\varphi &(\currule)\\
\end{plList}
\end{plContainer}

\end{proof-rule}


\begin{proof-rule}{la_generic}{veriT}
A step of the $\currule$ rule represents a tautological clause of linear
disequalities.  It can be checked by showing that the conjunction of
the negated disequalities is unsatisfiable. After the application of
some strengthening rules, the resulting conjunction is unsatisfiable,
even if integer variables are assumed to be real variables.

A linear inequality is of term of the form $\sum_{i=0}^{n}c_i\times{}t_i +
d_1\bowtie \sum_{i=n+1}^{m} c_i\times{}t_i + d_2$ where $\bowtie\;\in \{=, <,
>, \leq, \geq\}$, where $m\geq n$, $c_i, d_1, d_2$ are either integer or real
constants, and for each $i$ $c_i$ and $t_i$ have the same sort. We will write
$s_1 \bowtie s_2$.

Let $l_1,\dots, l_n$ be linear inequalities and $a_1, \dots, a_n$
rational numbers, then a $\currule$ step has the form

\begin{plListEq}
\proofsep& i.&
\varphi_1, \dots , \varphi_o
&(\currule; ; a_1, \dots, a_o)\\
\end{plListEq}

where $\varphi_i$ is either $\neg l_i$ or $l_i$, but never $s_1
\simeq s_2$.

If the current theory does not have rational numbers, then the $a_i$ are
printed using integer division. They should, nevertheless, be interpreted
as rational numbers. If $d_1$ or $d_2$ are $0$, they might not be printed.

To check the unsatisfiability of the negation of $\varphi_1\lor \dots
\lor \varphi_o$ one performs the following steps for each literal. For
each $i$, let $\varphi := \varphi_i$ and $a := a_i$.

\begin{enumerate}
    \item If $\varphi = s_1 > s_2$, then let $\varphi := s_1 \leq s_2$.
    	If $\varphi = s_1 \geq s_2$, then let $\varphi := s_1 < s_2$.
    	If $\varphi = s_1 < s_2$, then let $\varphi := s_1 \geq s_2$. 
    	If $\varphi = s_1 \leq s_2$, then let $\varphi := s_1 > s_2$. This negates
    	the literal.
   	\item If $\varphi = \neg (s_1 \bowtie s_2)$, then let $\varphi := s_1 \bowtie s_2$. 
    \item Replace $\varphi$ by $\sum_{i=0}^{n}c_i\times{}t_i - \sum_{i=n+1}^{m} c_i\times{}t_i
    \bowtie d$ where $d := d_2 - d_1$.
    \item Now $\varphi$ has the form $s_1 \bowtie d$. If all
    variables in $s_1$ are integer sorted: replace $\bowtie d$ according to
    table~\ref{tb:lageneric}.
    \item If $\bowtie$ is $\simeq$ replace $l$ by
    $\sum_{i=0}^{m}a\times{}c_i\times{}t_i \simeq a\times{}d$, otherwise replace it by
    $\sum_{i=0}^{m}|a|\times{}c_i\times{}t_i \simeq |a|\times{}d$.
\end{enumerate}

\begin{table}[h!]
\centering
\begin{tabular}{c|l l}
$\bowtie$ & If $d$ is an integer & Otherwise \\
\hline
$>$       & $\geq d + 1$ & $\geq \lfloor d\rfloor + 1$  \\
$\geq$    & $\geq d$     & $\geq \lfloor d\rfloor + 1$  \\
\end{tabular}
\caption{Strengthening rules for \texttt{la\_generic}.}
\label{tb:lageneric}
\end{table}

Finally, the sum of the resulting literals is trivially contradictory.
The sum
\begin{equation*}
    \sum_{k=1}^{o}\sum_{i=1}^{m^o}c_i^k*t_i^k \bowtie \sum_{k=1}^{o}d^k
\end{equation*}
where $c_i^k$ is the constant $c_i$ of literal $l_k$, $t_i^k$ is the term
$t_i$ of $l_k$, and $d^k$ is the constant $d$ of $l_k$. The operator
$\bowtie$ is $\simeq$ if all operators are $\simeq$, $>$ if all are
either $\simeq$ or $>$, and $\geq$ otherwise. The $a_i$ must be such
that the sum on the left-hand side is $0$ and the right-hand side is $>0$ (or
$\geq 0$ if $\bowtie is >$).

\begin{rule-example}
A simple $\currule$ step in the logic \texttt{LRA} might look like this:
\begin{minted}{smtlib2.py -x}
(step t10 (cl (not (> (f a) (f b))) (not (= (f a) (f b))))
    :rule la_generic :args (1.0 (- 1.0)))
\end{minted}
To verify this we have to check the insatisfiability of $f(a) > f(b) \land
f(a) = f(b)$ (Step 2). After step 3 we get $f(a) - f(b) > 0 \land
f(a) - f(b) = 0$. Since we don't have an integer sort in this logic step~4 does
not apply. Finally, after step~5 the conjunction is $f(a) - f(b) > 0 \land
- f(a) + f(b) = 0$. This sums to $0 > 0$, which is a contradiction.
\end{rule-example}

\begin{rule-example}
The following $\currule$ step is from a \texttt{QF_UFLIA} problem:
\begin{minted}{smtlib2.py -x}
(step t11 (cl (not (<= f3 0)) (<= (+ 1 (* 4 f3)) 1))
    :rule la_generic :args (1 (div 1 4)))
\end{minted}
After normalization we get $-f_3 \geq 0 \land 4\times f_3 > 0$.
This time step~4 applies and we can strengthen this to
$-f_3 \geq 0 \land 4\times f_3 \geq 1$ and after multiplication we get
$-f_3 \geq 0 \land f_3 \geq \frac{1}{4}$. Which sums to the contradiction
$\frac{1}{4} \geq 0$.
\end{rule-example}

\end{proof-rule}

\begin{proof-rule}{lia_generic}{veriT}
This rule is a placeholder rule for integer arithmetic solving. It takes the
same form as \texttt{la\_generic}, without the additional arguments.

\begin{plContainer}
\begin{plList}
\proofsep& i.&
\varphi_1, \dots , \varphi_n
&(\currule)\\
\end{plList}
\end{plContainer}
with $\varphi_i$ being linear inequalities. The disjunction
$\varphi_1\lor \dots \lor \varphi_n$ is a tautology in the theory of linear
integer arithmetic.

\paragraph{Remark.} Since this rule can introduce a disjunction of arbitrary
linear integer inequalities without any additional hints, proof checking
can be NP-hard. Hence, this rule should be avoided when possible.

\end{proof-rule}

\begin{proof-rule}{la_disequality}{veriT}
\begin{plContainer}
\begin{plList}
\proofsep& i.&
t_1 \simeq t_2 \lor \neg (t_1 \leq t_2) \lor \neg (t_2 \leq t_1)
&(\currule)\\
\end{plList}
\end{plContainer}
\end{proof-rule}

\begin{proof-rule}{la_totality}{veriT}
\begin{plContainer}
\begin{plList}
\proofsep& i.&
t_1 \leq t_2 \lor t_2 \leq t_1
&(\currule)\\
\end{plList}
\end{plContainer}
\end{proof-rule}

\begin{proof-rule}{la_tautology}{veriT}
This rule is a linear arithmetic tautology which can be checked without
sophisticated reasoning. It has either the form

\begin{plContainer}
\begin{plList}
\proofsep& i.&
\varphi
&(\currule)\\
\end{plList}
\end{plContainer}
where $\varphi$ is either a linear inequality $s_1 \bowtie s_2$
or $\neg(s_1 \bowtie s_2)$. After performing step 1 to 3 of the process for
checking the \proofRule{la\_generic} the result is trivially unsatisfiable.

The second form handles bounds on linear combinations. It is binary clause:
\begin{plContainer}
\begin{plList}
\proofsep& i.&
\varphi_1 \lor \varphi_2 % Checked: this is a proper disjunction not a cl
&(\currule)\\
\end{plList}
\end{plContainer}

It can be checked by using the procedure for \proofRule{la\_generic} with
while setting the arguments to $1$. Informally, the rule follows one of several
cases:
\begin{itemize}
    \item $\neg (s_1 \leq d_1) \lor s_1 \leq d_2$ where $d_1 \leq d_2$
    \item $s_1 \leq d_1 \lor \neg (s_1 \leq d_2)$ where $d_1 = d_2$
    \item $\neg (s_1 \geq d_1) \lor s_1 \geq d_2$ where $d_1 \geq d_2$
    \item $s_1 \geq d_1 \lor \neg (s_1 \geq d_2)$ where $d_1 = d_2$
    \item $\neg (s_1 \leq d_1) \lor \neg(s_1 \geq d_2)$ where $d_1 < d_2$
\end{itemize}
The inequalities $s_1 \bowtie d$ are are the result of applying normalization
as for the rule \proofRule{la\_generic}.

\end{proof-rule}

\begin{proof-rule}{bind}{veriT}
  The \currule{} rule is used to rename bound variables.

  \begin{plContainer}
    \begin{plSubList}
      \plLine\\
      \Gamma, y_1,\dots, y_n,  x_1 \mapsto y_1, \dots ,  x_n \mapsto y_n,
      \proofsep& j.& \varphi \leftrightarrow \varphi' &(\dots)\\
    \end{plSubList}
    \begin{plList}
      \Gamma \proofsep& k.&
      \forall x_1, \dots, x_n.\varphi \leftrightarrow \forall y_1, \dots, y_n. \varphi'  &(\currule)\\
    \end{plList}
  \end{plContainer}
  where the variables $y_1, \dots, y_n$ is not free in $\forall x_1,
  \dots, x_n.\varphi$.
\end{proof-rule}

\begin{proof-rule}{sko_ex}{veriT}
The \currule{} rule skolemizes existential quantifiers.

\begin{plContainer}
\begin{plSubList}
\plLine\\
\Gamma, x_1 \mapsto \varepsilon_1, \dots ,  x_n \mapsto \varepsilon_n
\proofsep& j.&\varphi \leftrightarrow \psi &(\dots)\\
\end{plSubList}
\begin{plList}
\Gamma \proofsep& k.&
\exists x_1, \dots, x_n.\varphi \leftrightarrow \psi  &(\currule)\\
\end{plList}
\end{plContainer}
where $\varepsilon_i$ stands for $\varepsilon x_i. (\exists x_{i+1}, \dots,
x_n. \varphi[x_1\mapsto \varepsilon_1,\dots,x_{i-1}\mapsto\varepsilon_{i-1}])$.
\end{proof-rule}

\begin{proof-rule}{sko_forall}{veriT}
The \currule{} rule skolemizes universal quantifiers.

\begin{plContainer}
\begin{center}
\parbox{\textwidth}{
\begin{plSubList}
\plLine\\
\Gamma, x_1 \mapsto (\varepsilon x_1.\neg\varphi), \dots,  x_n \mapsto (\varepsilon x_n.\neg\varphi)
\proofsep& j.&
\varphi \leftrightarrow \psi
&(\dots)\\
\end{plSubList}
\begin{plList}
\Gamma \proofsep& k.&
\forall x_1, \dots, x_n.\varphi \leftrightarrow \psi  &(\currule)\\
\end{plList}}
\end{center}
\end{plContainer}
\end{proof-rule}

\begin{proof-rule}{forall_inst}{veriT}
\begin{plContainer}
\begin{plList}
\proofsep& i.&
\neg (\forall x_1, \dots, x_n. P) \lor P[t_1/x_1]\dots[t_n/x_n]
&(\currule; ; (x_{k_1}, t_{k_1}), \dots, (x_{k_n}, t_{k_n}))\\
\end{plList}
\end{plContainer}
where $k_1, \dots, k_n$ is a permutation of $1, \dots, n$ and $x_i$ and
$k_i$ have the same sort. The arguments $(x_{k_i}, t_{k_i})$ are printed as
\smtinline{(:= xki tki)}.

\paragraph{Remark.} A rule simmilar to the \proofRule{let} rule would be
more appropriate.  The resulting proof would be more fine grained and this
would also be an opportunity to provide a proof for the clausification
as currently done by \proofRule{qnt_cnf}.

\end{proof-rule}

\begin{proof-rule}{refl}{veriT}
Either
\begin{plContainer}
\begin{plList}
\Gamma\proofsep& j.& t_1 \simeq t_2 &(\currule)\\
\end{plList}
\end{plContainer}
or
\begin{plContainer}
\begin{plList}
\Gamma\proofsep& j.& \varphi_1 \leftrightarrow \varphi_2 &(\currule)\\
\end{plList}
\end{plContainer}
\end{proof-rule}

where, if $\sigma = \operatorname{subst}(\Gamma)$,
the terms $\varphi_1\sigma$ and $\varphi_2$ (the formulas $P_1\sigma$ and $P_2$) are
syntactically equal up to the orientation of equalities.

\paragraph{Remark.} This is the only rule that requires the application of
the context.

\begin{proof-rule}{trans}{veriT}
Either
\begin{plContainer}
\begin{plList}
\Gamma\proofsep& i.& t_1 \simeq t_2 &(\dots)\\
\Gamma\proofsep& j.& t_2 \simeq t_3 &(\dots)\\
\Gamma\proofsep& k.& t_1 \simeq t_3 &(\currule; i, j)\\
\end{plList}
\end{plContainer}
or
\begin{plContainer}
\begin{plList}
\Gamma\proofsep& i.& \varphi_1 \leftrightarrow \varphi_2 &(\dots)\\
\Gamma\proofsep& j.& \varphi_2 \leftrightarrow \varphi_3 &(\dots)\\
\Gamma\proofsep& k.& \varphi_1 \leftrightarrow \varphi_3 &(\currule; i, j)\\
\end{plList}
\end{plContainer}

\begin{comment}{Mathias Fleury} The \proofRule{trans} rules comes in three
flavors that can be distinguished by the attribute: \begin{description}
\item[ordered and oriented] the equalities are given in the correct order
and are oriented in the right direction; \item[ordered and unoriented]
the equalities are given in the correct order, but
  the equalitis like ``\(t_1\simeq t_2\)'' can be used as ``\(t_1\simeq
  t_2\)'' or ``\(t_2\simeq t_1\)'';
\item[unordered and unoriented] the equalities are not ordered either.
\end{description}
\end{comment}

\end{proof-rule}

\begin{proof-rule}{cong}{veriT}
Either
\begin{plContainer}
\begin{plList}
\Gamma\proofsep& i_1.& t_1 \simeq u_1 &(\dots)\\
\Gamma\proofsep& i_n.& t_n \simeq u_n &(\dots)\\
\Gamma\proofsep& j.& \operatorname{f}(t_1, \dots, t_n) \simeq
\operatorname{f}(u_1, \dots, u_n) &(\currule;i_1, \dots, i_n)\\
\end{plList}
\end{plContainer}
where $\operatorname{f}$ is an $n$-ary function symbol,
or
\begin{plContainer}
\begin{plList}
\Gamma\proofsep& i_1.& \varphi_1 \simeq \psi_1 &(\dots)\\
\Gamma\proofsep& i_n.& \varphi_n \simeq \psi_n &(\dots)\\
\Gamma\proofsep& j.& \operatorname{P}(\varphi_1, \dots, \varphi_n) \leftrightarrow
\operatorname{P}(\psi_1, \dots, \psi_n) &(\currule;i_1, \dots, i_n)\\
\end{plList}
\end{plContainer}
where $\operatorname{P}$ is an $n$-ary predicate symbol.
\end{proof-rule}


\begin{proof-rule}{eq_reflexive}{veriT}
\begin{plContainer}
\begin{plList}
\proofsep& i.&
t \simeq t
&(\currule)\\
\end{plList}
\end{plContainer}
\end{proof-rule}

\begin{proof-rule}{eq_transitive}{veriT}
\begin{plContainer}
\begin{plList}
\proofsep& i.&
\neg (t_1 \simeq t_2) , \dots , \neg (t_{n-1} \simeq t_n) ,
t_1 \simeq t_n
&(\currule)\\
\end{plList}
\end{plContainer}
\end{proof-rule}

\begin{proof-rule}{eq_congruent}{veriT}
\begin{plContainer}
\begin{plList}
\proofsep& i.&
\neg (t_1 \simeq u_1) , \dots , \neg (t_n \simeq u_n) ,
f(t_1, \dots, t_n) \simeq f(u_1, \dots, u_n)
&(\currule)\\
\end{plList}
\end{plContainer}
\end{proof-rule}

\begin{proof-rule}{eq_congruent_pred}{veriT}
%\begin{compactproof}
%\begingroup\abovedisplayskip=0pt \belowdisplayskip=0pt%
\begin{plContainer}
\begin{plList}
\proofsep& i.&
\neg (t_1 \simeq u_1) , \dots , \neg (t_n \simeq u_n) ,
\neg (P(t_1, \dots, t_n) \simeq P(u_1, \dots, u_n))
&(\currule)\\
\end{plList}
\end{plContainer}
%\endgroup
%\end{compactproof}
%
\end{proof-rule}

\begin{proof-rule}{qnt_cnf}{veriT}
\begin{plContainer}
\begin{plList}
\proofsep& i.&
\neg(\forall x_1, \dots, x_n. \varphi) \lor \forall x_{k_1}, \dots, x_{k_m}.\varphi'
&(\currule)\\
\end{plList}
\end{plContainer}
This rule expresses clausification of a term under a universal
quantifier. This is used by conflicting instantiation. $\varphi'$ is one of the clause
of the clause normal form of $\varphi$. The variables $x_{k_1}, \dots, x_{k_m}$ are
a permutation of $x_1, \dots, x_n$ plus additional variables added by prenexing
$\varphi$. Normalization is performed in two phases. First, the negative normal form
is formed, then the result is prenexed. The result of the first step is $\Phi(\varphi, 1)$
where:

{\allowdisplaybreaks
  \begin{align*}
    \Phi(\neg \varphi, 1) &:= \Phi(\varphi, 0)\\
    \Phi(\neg \varphi, 0) &:= \Phi(\varphi, 1)\\
    \Phi(\varphi_1 \lor\dots\lor\varphi_n, 1) &:= \Phi(\varphi_1, 1)\lor\dots\lor\Phi(\varphi_n, 1)\\
    \Phi(\varphi_1 \land\dots\land\varphi_n, 1) &:= \Phi(\varphi_1, 1)\land\dots\land\Phi(\varphi_n, 1)\\
    \Phi(\varphi_1 \lor\dots\lor\varphi_n, 0) &:= \Phi(\varphi_1, 0)\land\dots\land\Phi(\varphi_n, 0)\\
    \Phi(\varphi_1 \land\dots\land\varphi_n, 0) &:= \Phi(\varphi_1, 0)\lor\dots\lor\Phi(\varphi_n, 0)\\
    \Phi(\varphi_1 \rightarrow \varphi_2, 1) &:= (\Phi(\varphi_1, 0) \lor \Phi(\varphi_2, 1)) \land
                                               (\Phi(\varphi_2, 0) \lor \Phi(\varphi_1, 1))\\
    \Phi(\varphi_1 \rightarrow \varphi_2, 0) &:= (\Phi(\varphi_1, 1) \land \Phi(\varphi_2, 0)) \lor
                                               (\Phi(\varphi_2, 1) \land \Phi(\varphi_1, 0))\\
    \Phi(\operatorname{ite} \varphi_1\;\varphi_2\;\varphi_3, 1) &:=
                                                                  (\Phi(\varphi_1, 0) \lor \Phi(\varphi_2, 1)) \land (\Phi(\varphi_1, 1) \lor \Phi(\varphi_3, 1))\\
    \Phi(\operatorname{ite} \varphi_1\;\varphi_2\;\varphi_3, 0) &:=
                                                                  (\Phi(\varphi_1, 1) \land \Phi(\varphi_2, 0)) \lor (\Phi(\varphi_1, 0) \land \Phi(\varphi_3, 0))\\
    \Phi(\forall x_1, \dots, x_n. \varphi, 1) &:= \forall x_1, \dots, x_n. \Phi(\varphi, 1)\\
    \Phi(\exists x_1, \dots, x_n. \varphi, 1) &:= \exists x_1, \dots, x_n. \Phi(\varphi, 1)\\
    \Phi(\forall x_1, \dots, x_n. \varphi, 0) &:= \exists x_1, \dots, x_n. \Phi(\varphi, 0)\\
    \Phi(\exists x_1, \dots, x_n. \varphi, 0) &:= \forall x_1, \dots, x_n. \Phi(\varphi, 0)\\
    \Phi(\varphi, 1) &:= \varphi\\
    \Phi(\varphi, 0) &:= \neg\varphi\\
  \end{align*}
}

\paragraph{Remark.} This is a placeholder rule that combines the many steps
done during clausification.

\end{proof-rule}

\begin{proof-rule}{and}{veriT}
  \begin{plContainer}
    \begin{plList}
      \proofsep& i.& \varphi_1 \land \dots \land \varphi_n&(\dots)\\
      \proofsep& j.& \varphi_i &(\currule; i)\\
    \end{plList}
  \end{plContainer}
\end{proof-rule}


\begin{proof-rule}{not_or}{veriT}
\begin{plContainer}
\begin{plList}
\proofsep& i.& \neg (\varphi_1 \lor \dots \lor \varphi_n) &(\dots)\\
\proofsep& j.& \neg \varphi_i &(\currule; i)\\
\end{plList}
\end{plContainer}
\end{proof-rule}

\begin{proof-rule}{or}{veriT}
\begin{plContainer}
\begin{plList}
\proofsep& i.& \varphi_1 \lor \dots \lor \varphi_n &(\dots)\\
\proofsep& j.& \varphi_1 , \dots , \varphi_n &(\currule; i)\\
\end{plList}
\end{plContainer}

\paragraph{Remark.} This rule deconstructs the \smtinline{or} operator
into a clause denoted by $\smtinline{cl}$.

\begin{rule-example}
An application of the \currule{} rule.

\begin{minted}{smtlib2.py -x}
(step t15 (cl (or (= a b) (not (<= a b)) (not (<= b a))))
    :rule la_disequality)
(step t16 (cl     (= a b) (not (<= a b)) (not (<= b a)))
    :rule or :premises (t15))
\end{minted}
\end{rule-example}
\end{proof-rule}

\begin{proof-rule}{not_and}{veriT}
\begin{plContainer}
\begin{plList}
\proofsep& i.& \neg (\varphi_1 \land \dots \land \varphi_n) &(\dots)\\
\proofsep& j.& \neg\varphi_1 , \dots , \neg\varphi_n &(\currule; i)\\
\end{plList}
\end{plContainer}

\end{proof-rule}

\begin{proof-rule}{xor1}{veriT}
\begin{plContainer}
\begin{plList}
\proofsep& i.& \operatorname{xor} \varphi_1\;\varphi_2 &(\dots)\\
\proofsep& j.& \varphi_1 , \varphi_2 &(\currule; i)\\
\end{plList}
\end{plContainer}
\end{proof-rule}

\begin{proof-rule}{xor2}{veriT}
\begin{plContainer}
\begin{plList}
\proofsep& i.& \operatorname{xor} \varphi_1\;\varphi_2 &(\dots)\\
\proofsep& j.& \neg\varphi_1 , \neg\varphi_2 &(\currule; i)\\
\end{plList}
\end{plContainer}
\end{proof-rule}

\begin{proof-rule}{not_xor1}{veriT}
\begin{plContainer}
\begin{plList}
\proofsep& i.& \neg(\operatorname{xor} \varphi_1\;\varphi_2) &(\dots)\\
\proofsep& j.& \varphi_1 , \neg\varphi_2 &(\currule; i)\\
\end{plList}
\end{plContainer}
\end{proof-rule}

\begin{proof-rule}{not_xor2}{veriT}
\begin{plContainer}
\begin{plList}
\proofsep& i.& \neg(\operatorname{xor} \varphi_1\;\varphi_2) &(\dots)\\
\proofsep& j.& \neg\varphi_1 , \varphi_2 &(\currule; i)\\
\end{plList}
\end{plContainer}
\end{proof-rule}

\begin{proof-rule}{implies}{veriT}
\begin{plContainer}
\begin{plList}
\proofsep& i.& \varphi_1\rightarrow\varphi_2 &(\dots)\\
\proofsep& j.& \neg\varphi_1 , \varphi_2 &(\currule; i)\\
\end{plList}
\end{plContainer}
\end{proof-rule}

\begin{proof-rule}{not_implies1}{veriT}
\begin{plContainer}
\begin{plList}
\proofsep& i.& \neg (\varphi_1\rightarrow\varphi_2) &(\dots)\\
\proofsep& j.& \varphi_1 &(\currule; i)\\
\end{plList}
\end{plContainer}
\end{proof-rule}

\begin{proof-rule}{not_implies2}{veriT}
\begin{plContainer}
\begin{plList}
\proofsep& i.& \neg (\varphi_1\rightarrow\varphi_2) &(\dots)\\
\proofsep& j.& \neg\varphi_2 &(\currule; i)\\
\end{plList}
\end{plContainer}
\end{proof-rule}

\begin{proof-rule}{equiv1}{veriT}
\begin{plContainer}
\begin{plList}
\proofsep& i.& \varphi_1\leftrightarrow\varphi_2 &(\dots)\\
\proofsep& j.& \neg\varphi_1 , \varphi_2 &(\currule; i)\\
\end{plList}
\end{plContainer}
\end{proof-rule}

\begin{proof-rule}{equiv2}{veriT}
\begin{plContainer}
\begin{plList}
\proofsep& i.& \varphi_1\leftrightarrow\varphi_2 &(\dots)\\
\proofsep& j.& \varphi_1 , \neg\varphi_2 &(\currule; i)\\
\end{plList}
\end{plContainer}
\end{proof-rule}

\begin{proof-rule}{not_equiv1}{veriT}
\begin{plContainer}
\begin{plList}
\proofsep& i.& \neg(\varphi_1\leftrightarrow\varphi_2) &(\dots)\\
\proofsep& j.& \varphi_1 , \varphi_2 &(\currule; i)\\
\end{plList}
\end{plContainer}
\end{proof-rule}

\begin{proof-rule}{not_equiv2}{veriT}
\begin{plContainer}
\begin{plList}
\proofsep& i.& \neg(\varphi_1\leftrightarrow\varphi_2) &(\dots)\\
\proofsep& j.& \neg\varphi_1 , \neg\varphi_2 &(\currule; i)\\
\end{plList}
\end{plContainer}
\end{proof-rule}

\begin{proof-rule}{and_pos}{veriT}
\begin{plContainer}
\begin{plList}
\proofsep& i.& \neg (\varphi_1 \land \dots \land \varphi_n) , \varphi_k &(\currule)\\
\end{plList}
\end{plContainer}
with $1 \leq k \leq n$.
\end{proof-rule}

\begin{proof-rule}{and_neg}{veriT}
\begin{plContainer}
\begin{plList}
\proofsep& i.& (\varphi_1 \land \dots \land \varphi_n), \neg\varphi_1
	, \dots , \neg\varphi_n &(\currule)\\
\end{plList}
\end{plContainer}
\end{proof-rule}

\begin{proof-rule}{or_pos}{veriT}
\begin{plContainer}
\begin{plList}
\proofsep& i.& \neg (\varphi_1 \lor \dots \lor \varphi_n) , \varphi_1 , \dots
	, \varphi_n &(\currule)\\
\end{plList}
\end{plContainer}
\end{proof-rule}

\begin{proof-rule}{or_neg}{veriT}
\begin{plContainer}
\begin{plList}
\proofsep& i.& (\varphi_1 \lor \dots \lor \varphi_n) , \neg \varphi_k &(\currule)\\
\end{plList}
\end{plContainer}
with $1 \leq k \leq n$.
\end{proof-rule}

\begin{proof-rule}{xor_pos1}{veriT}
\begin{plContainer}
\begin{plList}
\proofsep& i.& \neg (\varphi_1 \operatorname{xor} \varphi_2) , \varphi_1 , \varphi_2 &(\currule)\\
\end{plList}
\end{plContainer}
\end{proof-rule}

\begin{proof-rule}{xor_pos2}{veriT}
\begin{plContainer}
\begin{plList}
\proofsep& i.& \neg (\varphi_1 \operatorname{xor} \varphi_2)
, \neg \varphi_1, \neg \varphi_2 &(\currule)\\
\end{plList}
\end{plContainer}
\end{proof-rule}

\begin{proof-rule}{xor_neg1}{veriT}
\begin{plContainer}
\begin{plList}
\proofsep& i.& \varphi_1 \operatorname{xor} \varphi_2
, \varphi_1 , \neg \varphi_2 &(\currule)\\
\end{plList}
\end{plContainer}
\end{proof-rule}

\begin{proof-rule}{xor_neg2}{veriT}
\begin{plContainer}
\begin{plList}
\proofsep& i.& \varphi_1 \operatorname{xor} \varphi_2
, \neg \varphi_1 , \varphi_2 &(\currule)\\
\end{plList}
\end{plContainer}
\end{proof-rule}

\begin{proof-rule}{implies_pos}{veriT}
\begin{plContainer}
\begin{plList}
\proofsep& i.& \neg (\varphi_1 \rightarrow \varphi_2)
, \neg \varphi_1 , \varphi_2 &(\currule)\\
\end{plList}
\end{plContainer}
\end{proof-rule}

\begin{proof-rule}{implies_neg1}{veriT}
\begin{plContainer}
\begin{plList}
\proofsep& i.& \varphi_1 \rightarrow \varphi_2
, \varphi_1 &(\currule)\\
\end{plList}
\end{plContainer}
\end{proof-rule}

\begin{proof-rule}{implies_neg2}{veriT}
\begin{plContainer}
\begin{plList}
\proofsep& i.& \varphi_1 \rightarrow \varphi_2
, \neg \varphi_2 &(\currule)\\
\end{plList}
\end{plContainer}
\end{proof-rule}

\begin{proof-rule}{equiv_pos1}{veriT}
\begin{plContainer}
\begin{plList}
\proofsep& i.&
\neg (\varphi_1 \leftrightarrow \varphi_2) , \varphi_1 , \neg \varphi_2
&(\currule)\\
\end{plList}
\end{plContainer}
\end{proof-rule}

\begin{proof-rule}{equiv_pos2}{veriT}
\begin{plContainer}
\begin{plList}
\proofsep& i.&
\neg (\varphi_1 \leftrightarrow \varphi_2) , \neg \varphi_1 , \varphi_2
&(\currule)\\
\end{plList}
\end{plContainer}
\end{proof-rule}

\begin{proof-rule}{equiv_neg1}{veriT}
\begin{plContainer}
\begin{plList}
\proofsep& i.&
\varphi_1 \leftrightarrow \varphi_2 , \neg \varphi_1 , \neg \varphi_2
&(\currule)\\
\end{plList}
\end{plContainer}
\end{proof-rule}

\begin{proof-rule}{equiv_neg2}{veriT}
\begin{plContainer}
\begin{plList}
\proofsep& i.&
\varphi_1 \leftrightarrow \varphi_2 , \varphi_1 , \varphi_2
&(\currule)\\
\end{plList}
\end{plContainer}
\end{proof-rule}

\begin{proof-rule}{ite1}{veriT}
\begin{plContainer}
\begin{plList}
\proofsep& i.& \operatorname{ite}\varphi_1\;\varphi_2\;\varphi_3 &(\dots)\\
\proofsep& j.& \varphi_1 , \varphi_3 &(\currule; i)\\
\end{plList}
\end{plContainer}
\end{proof-rule}

\begin{proof-rule}{ite2}{veriT}
\begin{plContainer}
\begin{plList}
\proofsep& i.& \operatorname{ite}\varphi_1\;\varphi_2\;\varphi_3 &(\dots)\\
\proofsep& j.& \neg\varphi_1 , \varphi_2 &(\currule; i)\\
\end{plList}
\end{plContainer}
\end{proof-rule}

\begin{proof-rule}{ite_pos1}{veriT}
  \begin{plContainer}
    \begin{plList}
      \proofsep& i.&
      \neg (\operatorname{ite} \varphi_1\;\varphi_2\;\varphi_3) , \varphi_1 , \varphi_3
      &(\currule)\\
    \end{plList}
  \end{plContainer}
\end{proof-rule}

\begin{proof-rule}{ite_pos2}{veriT}
  \begin{plContainer}
    \begin{plList}
      \proofsep& i.&
      \neg (\operatorname{ite} \varphi_1\;\varphi_2\;\varphi_3) , \neg \varphi_1 , \varphi_2
      &(\currule)\\
    \end{plList}
  \end{plContainer}
\end{proof-rule}

\begin{proof-rule}{ite_neg1}{veriT}
  \begin{plContainer}
    \begin{plList}
      \proofsep& i.&
      \operatorname{ite} \varphi_1\;\varphi_2\;\varphi_3 , \varphi_1 , \neg \varphi_3
      &(\currule)\\
    \end{plList}
  \end{plContainer}
\end{proof-rule}

\begin{proof-rule}{ite_neg2}{veriT}
  \begin{plContainer}
    \begin{plList}
      \proofsep& i.&
      \operatorname{ite} \varphi_1\;\varphi_2\;\varphi_3 , \neg \varphi_1 , \neg \varphi_2
      &(\currule)\\
    \end{plList}
  \end{plContainer}
\end{proof-rule}

\begin{proof-rule}{not_ite1}{veriT}
\begin{plContainer}
\begin{plList}
\proofsep& i.& \neg(\operatorname{ite}\varphi_1\;\varphi_2\;\varphi_3) &(\dots)\\
\proofsep& j.& \varphi_1 , \neg\varphi_3 &(\currule; i)\\
\end{plList}
\end{plContainer}
\end{proof-rule}

\begin{proof-rule}{not_ite2}{veriT}
\begin{plContainer}
\begin{plList}
\proofsep& i.& \neg(\operatorname{ite}\varphi_1\;\varphi_2\;\varphi_3) &(\dots)\\
\proofsep& j.& \neg\varphi_1 , \neg\varphi_2 &(\currule; i)\\
\end{plList}
\end{plContainer}
\end{proof-rule}

\begin{proof-rule}{connective_def}{veriT}
  This rule is used to replace connectives by their definition. It can be one
  of the following:

  \begin{plContainer}
    \begin{plList}
      \Gamma\proofsep&
      i.& \varphi_1\operatorname{xor}\varphi_2 \leftrightarrow
      (\neg\varphi_1 \land \varphi_2) \lor (\varphi_1 \land \neg\varphi_2)
      &(\currule)\\
    \end{plList}
  \end{plContainer}

  \begin{plContainer}
    \begin{plList}
      \Gamma\proofsep&
      i.& \varphi_1\leftrightarrow\varphi_2 \leftrightarrow
      (\varphi_1 \rightarrow \varphi_2) \land (\varphi_2 \rightarrow \varphi_1)
      &(\currule)\\
    \end{plList}
  \end{plContainer}

  \begin{plContainer}
    \begin{plList}
      \Gamma\proofsep&
      i.& \operatorname{ite} \varphi_1\;\varphi_2\;\varphi_3 \leftrightarrow
      (\varphi_1 \rightarrow \varphi_2) \land (\neg\varphi_1 \rightarrow \varphi_3)
      &(\currule)\\
    \end{plList}
  \end{plContainer}

  \begin{plContainer}
    \begin{plList}
      \Gamma\proofsep&
      i.& \operatorname{\forall x_1, \dots, x_n.} \varphi \leftrightarrow \operatorname{\neg\exists x_1, \dots, x_n.} \neg\varphi
      &(\currule)\\
    \end{plList}
  \end{plContainer}

\end{proof-rule}

\begin{proof-rule}{and_simplify}{veriT}
This rule simplifies an $\land$ term by applying equivalent
transformations as long as possible. Hence, the general form is

\begin{plContainer}
\begin{plList}
\Gamma\proofsep&
i.& \varphi_1\land \dots\land\varphi_n \leftrightarrow \psi
&(\currule)\\
\end{plList}
\end{plContainer}
where $\psi$ is the transformed term.

The possible transformations are:
\begin{itemize}
    \item $\top \land \dots \land \top \leftrightarrow \top$
    \item $\varphi_1 \land \dots \land \varphi_n \leftrightarrow \varphi_1
    \land \dots \land \varphi_{n'} $ where the right hand side has all
    $\top$ literals removed.
    \item $\varphi_1 \land \dots \land \varphi_n \leftrightarrow \varphi_1
    \land \dots \land \varphi_{n'} $ where the right hand side has all
    repeated literals removed.
    \item $\varphi_1 \land\dots\land \bot\land\dots \land \varphi_n \leftrightarrow \bot$
    \item $\varphi_1 \land\dots\land \varphi_i\land \dots \land \varphi_j\land\dots \land \varphi_n \leftrightarrow \bot$ if
      $\varphi_i = \bar{\varphi_j}$
\end{itemize}
\end{proof-rule}

\begin{proof-rule}{or_simplify}{veriT}
This rule simplifies an $\lor$ term by applying equivalent
transformations as long as possible. Hence, the general form is

\begin{plContainer}
\begin{plList}
\Gamma\proofsep&
i.& (\varphi_1\lor \dots\lor\varphi_n) \leftrightarrow \psi
&(\currule)\\
\end{plList}
\end{plContainer}
where $\psi$ is the transformed term.

The possible transformations are:
\begin{itemize}
    \item $\bot \lor \dots \lor \bot \leftrightarrow \bot$
    \item $\varphi_1 \lor \dots \lor \varphi_n \leftrightarrow \varphi_1
    \lor \dots \lor \varphi_{n'} $ where the right hand side has all
    $\bot$ literals removed.
    \item $\varphi_1 \lor \dots \lor \varphi_n \leftrightarrow \varphi_1
    \lor \dots \lor \varphi_{n'} $ where the right hand side has all
    repeated literals removed.
    \item $\varphi_1 \lor\dots\lor \top\lor\dots \lor \varphi_n \leftrightarrow \top$
    \item $\varphi_1 \lor\dots\lor \varphi_i\lor \dots \lor \varphi_j\lor\dots \lor \varphi_n \leftrightarrow \top$ if
    	$\varphi_i = \bar{\varphi_j}$
\end{itemize}
\end{proof-rule}

\begin{proof-rule}{not_simplify}{veriT}
This rule simplifies an $\neg$ term by applying equivalent
transformations as long as possible. Hence, the general form is

\begin{plContainer}
\begin{plList}
\Gamma\proofsep&
i.& \neg\varphi \leftrightarrow \psi
&(\currule)\\
\end{plList}
\end{plContainer}
where $\psi$ is the transformed term.

The possible transformations are:
\begin{itemize}
    \item $\neg (\neg \varphi) \leftrightarrow \varphi$
    \item $\neg \bot \leftrightarrow \top$
    \item $\neg \top \leftrightarrow \bot$
\end{itemize}
\end{proof-rule}

\begin{proof-rule}{implies_simplify}{veriT}
This rule simplifies an $\rightarrow$ term by applying equivalent
transformations as long as possible. Hence, the general form is

\begin{plContainer}
\begin{plList}
\Gamma\proofsep&
i.& \varphi_1\rightarrow \varphi_2 \leftrightarrow \psi
&(\currule)\\
\end{plList}
\end{plContainer}
where $\psi$ is the transformed term.

The possible transformations are:
\begin{itemize}
    \item $\neg \varphi_1 \rightarrow \neg \varphi_2 \leftrightarrow  \varphi_2\rightarrow \varphi_1$
    \item $\bot \rightarrow  \varphi \leftrightarrow \top$
    \item $ \varphi \rightarrow \top \leftrightarrow \top$
    \item $\top \rightarrow  \varphi \leftrightarrow  \varphi$
		\item $ \varphi \rightarrow \bot \leftrightarrow \neg \varphi$
		\item $ \varphi \rightarrow  \varphi \leftrightarrow \top$
		\item $\neg \varphi \rightarrow  \varphi \leftrightarrow  \varphi$
		\item $ \varphi \rightarrow \neg \varphi \leftrightarrow \neg \varphi$
		\item $( \varphi_1\rightarrow \varphi_2)\rightarrow \varphi_2 \leftrightarrow  \varphi_1\lor \varphi_2$
\end{itemize}
\end{proof-rule}

\begin{proof-rule}{equiv_simplify}{veriT}
This rule simplifies an $\leftrightarrow$ term by applying equivalent
transformations as long as possible. Hence, the general form is

\begin{plContainer}
\begin{plList}
\Gamma\proofsep&
i.&  (\varphi_1\leftrightarrow \varphi_2) \leftrightarrow \psi
&(\currule)\\
\end{plList}
\end{plContainer}
where $\psi$ is the transformed term.

The possible transformations are:
\begin{itemize}
    \item $(\neg \varphi_1 \leftrightarrow \neg \varphi_2) \leftrightarrow ( \varphi_1\leftrightarrow \varphi_2)$
    \item $( \varphi\leftrightarrow \varphi) \leftrightarrow \top$
    \item $( \varphi\leftrightarrow \neg \varphi) \leftrightarrow \bot$
    \item $(\neg \varphi\leftrightarrow  \varphi) \leftrightarrow \bot$
    \item $(\top \leftrightarrow  \varphi) \leftrightarrow  \varphi$
    \item $( \varphi \leftrightarrow \top) \leftrightarrow  \varphi$
    \item $(\bot \leftrightarrow  \varphi) \leftrightarrow \neg \varphi$
    \item $( \varphi \leftrightarrow \bot) \leftrightarrow \neg \varphi$
\end{itemize}
\end{proof-rule}

\begin{proof-rule}{bool_simplify}{veriT}
This rule simplifies a boolean term by applying equivalent
transformations as long as possible. Hence, the general form is

\begin{plContainer}
\begin{plList}
\Gamma\proofsep&
i.& \varphi\leftrightarrow \psi
&(\currule)\\
\end{plList}
\end{plContainer}
where $\psi$ is the transformed term.

The possible transformations are:
\begin{itemize}
	\item $\neg(\varphi_1\rightarrow \varphi_2) \leftrightarrow (\varphi_1 \land \neg \varphi_2)$
	\item $\neg(\varphi_1\lor \varphi_2) \leftrightarrow (\neg \varphi_1 \land \neg \varphi_2)$
	\item $\neg(\varphi_1\land \varphi_2) \leftrightarrow (\neg \varphi_1 \lor \neg \varphi_2)$
	\item $(\varphi_1 \rightarrow (\varphi_2\rightarrow \varphi_3)) \leftrightarrow (\varphi_1\land \varphi_2) \rightarrow \varphi_3$
	\item $((\varphi_1\rightarrow \varphi_2)\rightarrow \varphi_2)  \leftrightarrow (\varphi_1\lor \varphi_2)$
	\item $(\varphi_1 \land (\varphi_1\rightarrow \varphi_2)) \leftrightarrow (\varphi_1 \land \varphi_2)$
	\item $((\varphi_1\rightarrow \varphi_2) \land \varphi_1) \leftrightarrow (\varphi_1 \land \varphi_2)$
\end{itemize}
\end{proof-rule}


\begin{proof-rule}{ac_simp}{veriT}
  This rule simplifies nested occurences of $\lor$ or $\land$:

  \begin{plContainer}
    \begin{plList}
      \Gamma\proofsep&
      i.& \psi  \leftrightarrow \varphi_1 \circ\dots\circ\varphi_n
      &(\currule)\\
    \end{plList}
  \end{plContainer}
  where $\circ \in \{\lor, \land\}$ and $\psi$ is a nested application of $\circ$.
  The literals $\varphi_i$ are literals of the flattening of $\psi$ with duplicates
  removed.
\end{proof-rule}

\subsubsection*{If-Then-Else Operators}
\label{sec:ite-simp-rules}

\begin{proof-rule}{ite_simplify}{veriT}
  This rule simplifies an if-then-else term by applying equivalent
  transformations until fix point\footnote{Note however that the order of the
    application is important, since the set of rules is not confluent. For
    example, the term $(\operatorname{ite} \top \; t_1 \; t_2 \leftrightarrow
    t_1)$ can be simplified into both $p$ and $(not (not p))$ depending on the
    order of applications.}
  %
Depending on the sort of the
$\operatorname{ite}$-term the rule can have one of two forms. If the sort
is $\mathbf{Bool}$ it has the form
\begin{plContainer}
\begin{plList}
\Gamma\proofsep&
i.& \operatorname{ite} \varphi\;t_1\;t_2 \leftrightarrow \psi
&(\currule)\\
\end{plList}
\end{plContainer}
where $\psi$ is the transformed term.

Otherwise, it has the form
\begin{plContainer}
\begin{plList}
\Gamma\proofsep&
i.& \operatorname{ite} \varphi\;t_1\;t_2 \simeq u
&(\currule)\\
\end{plList}
\end{plContainer}
where $u$ is the transformed term.

The possible transformations are:
\begin{itemize}
    \item $\operatorname{ite} \top      \; t_1 \; t_2 \leftrightarrow t_1$
    \item $\operatorname{ite} \bot      \; t_1 \; t_2 \leftrightarrow t_2$
    \item $\operatorname{ite} \psi      \; t \; t \leftrightarrow t$
    \item $\operatorname{ite} \neg \varphi \; t_1 \; t_2 \leftrightarrow \operatorname{ite} \varphi \; t_2 \; t_1$
    \item $\operatorname{ite} \psi \; (\operatorname{ite} \psi\;t_1\;t_2)\; t_3 \leftrightarrow
			\operatorname{ite} \psi\; t_1\; t_3$
    \item $\operatorname{ite} \psi \; t_1\; (\operatorname{ite} \psi\;t_2\;t_3) \leftrightarrow
			\operatorname{ite} \psi\; t_1\; t_3$
    \item $\operatorname{ite} \psi \; \top\; \bot \leftrightarrow \psi$
    \item $\operatorname{ite} \psi \; \bot\; \top \leftrightarrow \neg\psi$
    \item $\operatorname{ite} \psi \; \top \; \varphi \leftrightarrow \psi\lor\varphi$
    \item $\operatorname{ite} \psi \; \varphi\;\bot \leftrightarrow \psi\land\varphi$
    \item $\operatorname{ite} \psi \; \bot\; \varphi \leftrightarrow \neg\psi\land\varphi$
    \item $\operatorname{ite} \psi \; \varphi\;\top \leftrightarrow \neg\psi\lor\varphi$
\end{itemize}
\end{proof-rule}

\begin{proof-rule}{qnt_simplify}{veriT}
  This rule simplifies a $\forall$-formula with a constant predicate.

  \begin{plContainer}
    \begin{plList}
      \Gamma\proofsep&
      i.&  \forall x_1, \dots, x_n. \varphi \leftrightarrow \varphi
      &(\currule)\\
    \end{plList}
  \end{plContainer}
  where $\varphi$ is either $\top$ or $\bot$.
\end{proof-rule}

\begin{proof-rule}{onepoint}{veriT}
The {\currule} rule is the ``one-point-rule''. That is: it eliminates quantified
variables that can only have one value.

\begin{plContainer}
\begin{plContainer}
\begin{plSubList}
\plLine\\
\Gamma, x_{k_1},\dots, x_{k_m},  x_{j_1} \mapsto t_{j_1}, \dots ,  x_{j_o} \mapsto t_{j_o},
\proofsep& j.& \varphi \leftrightarrow \varphi' &(\dots)\\
\end{plSubList}
\begin{plList}
\Gamma \proofsep& k.&
Q x_1, \dots, x_n.\varphi \leftrightarrow Q x_{k_1}, \dots, x_{k_m}. \varphi'  &(\currule)\\
\end{plList}
\end{plContainer}
\end{plContainer}
where $Q\in\{\forall, \exists\}$,  $n = m + o$,  $k_1, \dots, k_m$ and
$j_1, \dots, j_o$ are monotone
mappings to $1, \dots, n$, and no $x_{k_i}$ appears in $x_{j_1}, \dots, x_{j_o}$.

The terms $t_{j_1}, \dots, t_{j_o}$ are the points of the variables
$x_{j_1}, \dots, x_{j_o}$. Points are defined by equalities $x_i\simeq t_i$
with positive polarity in the term $\varphi$.

\paragraph{Remark.} Since an eliminated variable $x_i$ might appear free in a
term $t_j$, it is necessary to replace $x_i$ with $t_i$ inside $t_j$. While
this substitution is performed correctly, the proof for it is currently
missing.

\begin{rule-example}
An application of the $\currule$ rule on the term $\forall x, y.\, x \simeq y
\rightarrow f(x)\land f(y)$ look like this:

\begin{minted}{smtlib2.py -x}
(anchor :step t3 :args ((:= y x)))
(step t3.t1 (cl (= x y)) :rule refl)
(step t3.t2 (cl (= (= x y) (= x x)))
    :rule cong :premises (t3.t1))
(step t3.t3 (cl (= x y)) :rule refl)
(step t3.t4 (cl (= (f y) (f x)))
    :rule cong :premises (t3.t3))
(step t3.t5 (cl (= (and (f x) (f y)) (and (f x) (f x))))
    :rule cong :premises (t3.t4))
(step t3.t6 (cl (=
        (=> (= x y) (and (f x) (f y)))
        (=> (= x x) (and (f x) (f x)))))
    :rule cong :premises (t3.t2 t3.t5))
(step t3 (cl (=
        (forall ((x S) (y S)) (=> (= x y) (and (f x) (f y))))
        (forall ((x S)) (=> (= x x) (and (f x) (f x))))))
    :rule qnt_simplify)
\end{minted}
\end{rule-example}

\end{proof-rule}

\begin{proof-rule}{qnt_join}{veriT}
  \begin{plContainer}
    \begin{plList}
      \Gamma\proofsep& i.&
      Q x_1, \dots, x_n. Q x_{n+1}, \dots, x_{m}. \varphi
      \leftrightarrow
      Q x_{k_1}, \dots, x_{k_o}. \varphi
      &(\currule)\\
    \end{plList}
  \end{plContainer}
  where $m > n$ and $Q\in\{\forall, \exists\}$. Furthermore, $k_1, \dots, k_o$ is a monotonic
  map to $1, \dots, m$ such that $x_{k_1}, \dots, x_{k_o}$ are pairwise
  distinct, and $\{x_1, \dots, x_m\} = \{x_{k_1}, \dots, x_{k_o}\}$.
\end{proof-rule}

\begin{proof-rule}{qnt_rm_unused}{veriT}
  \begin{plContainer}
    \begin{plList}
      \Gamma\proofsep& i.&
      Q x_1, \dots, x_n. \varphi \leftrightarrow Q x_{k_1}, \dots, x_{k_m}.\varphi
      &(\currule)\\
    \end{plList}
  \end{plContainer}
  where $m \leq n$ and $Q\in\{\forall, \exists\}$. Furthermore, $k_1, \dots, k_m$ is
  a monotonic map to $1, \dots, n$ and if $x\in \{x_j\; |\; j \in \{1, \dots,
  n\} \land j\not\in \{k_1, \dots, k_m\}\}$ then $x$ is not free in $P$.

\end{proof-rule}

\begin{proof-rule}{eq_simplify}{veriT}
  This rule simplifies an $\simeq$ term by applying equivalent
  transformations as long as possible. Hence, the general form is

  \begin{plContainer}
    \begin{plList}
      \Gamma\proofsep&
      i.& t_1\simeq t_2 \leftrightarrow \varphi
      &(\currule)\\
    \end{plList}
  \end{plContainer}
  where $\psi$ is the transformed term.

  The possible transformations are:
  \begin{itemize}
  \item $t \simeq t \leftrightarrow \top$
  \item $t_1 \simeq t_2 \leftrightarrow \bot$ if $t_1$ and $t_2$ are different numeric constants.
  \item $\neg (t \simeq t) \leftrightarrow \bot$ if $t$ is a numeric constant.
  \end{itemize}
\end{proof-rule}

\begin{proof-rule}{div_simplify}{veriT}
This rule simplifies a division by applying equivalent
transformations. The general form is

\begin{plContainer}
\begin{plList}
\Gamma\proofsep&
i.& (t_1 / t_2) \simeq t_3
&(\currule)\\
\end{plList}
\end{plContainer}

\noindent The possible transformations are:
\begin{itemize}
  \item $t/t = 1$
  \item $t/1 = t$
	\item $t_1 / t_2 = t_3$
		if $t_1$ and $t_2$ are constants and $t_3$ is $t_1$
		divided by $t_2$ according to the semantic of the current theory.
\end{itemize}
\end{proof-rule}

\begin{proof-rule}{prod_simplify}{veriT}
This rule simplifies a product by applying equivalent
transformations as long as possible. The general form is

\begin{plContainer}
\begin{plList}
\Gamma\proofsep&
i.& t_1\times\dots\times t_n \simeq u
&(\currule)\\
\end{plList}
\end{plContainer}
where $u$ is either a constant or a product.

The possible transformations are:
\begin{itemize}
    \item $t_1\times\dots\times t_n = u$ where all
    $t_i$ are constants and $u$ is their product.
    \item $t_1\times\dots\times t_n = 0$ if any
    $t_i$ is $0$.
    \item $t_1\times\dots\times t_n =
			c \times t_{k_1}\times\dots\times t_{k_n}$ where $c$
			ist the product of the constants of $t_1, \dots, t_n$ and
			$t_{k_1}, \dots, t_{k_n}$ is $t_1, \dots, t_n$
			with the constants removed.
    \item $t_1\times\dots\times t_n =
			t_{k_1}\times\dots\times t_{k_n}$: same as above if $c$ is
			$1$.
\end{itemize}
\end{proof-rule}

\begin{proof-rule}{unary_minus_simplify}{veriT}
This rule is either

\begin{plContainer}
\begin{plList}
\Gamma\proofsep&
i.& - (-t) \simeq t &(\currule)\\
\end{plList}
\end{plContainer}
or
\begin{plContainer}
\begin{plList}
\Gamma\proofsep&
i.& -t \simeq u &(\currule)\\
\end{plList}
\end{plContainer}
where $u$ is the negated numerical constant $t$.
\end{proof-rule}

\begin{proof-rule}{minus_simplify}{veriT}
This rule simplifies a subtraction by applying equivalent
transformations. The general form is

\begin{plContainer}
\begin{plList}
\Gamma\proofsep&
i.& t_1 - t_2 \simeq u &(\currule)\\
\end{plList}
\end{plContainer}

\noindent The possible transformations are:
\begin{itemize}
    \item $t - t = 0$
    \item $t_1 - t_2 = t_3$ where $t_1$
    and $t_2$ are numerical constants and $t_3$ is $t_2$ subtracted
    from~$t_1$.
    \item $t - 0 = t$
    \item $0 - t = -t$
\end{itemize}
\end{proof-rule}

\begin{proof-rule}{sum_simplify}{veriT}
This rule simplifies a sum by applying equivalent
transformations as long as possible. The general form is

\begin{plContainer}
\begin{plList}
\Gamma\proofsep&
i.& t_1+\dots+t_n \simeq u &(\currule)\\
\end{plList}
\end{plContainer}
where $u$ is either a constant or a product.

The possible transformations are:
\begin{itemize}
    \item $t_1+\dots+t_n = c$ where all
    $t_i$ are constants and $c$ is their sum.
    \item $t_1+\dots+t_n =
			c + t_{k_1}+\dots+t_{k_n}$ where $c$
			ist the sum of the constants of $t_1, \dots, t_n$ and
			$t_{k_1}, \dots, t_{k_n}$ is $t_1, \dots, t_n$
			with the constants removed.
    \item $t_1+\dots+t_n =
			t_{k_1}+\dots+t_{k_n}$: same as above if $c$ is
			$0$.
\end{itemize}
\end{proof-rule}

\begin{proof-rule}{comp_simplify}{veriT}
This rule simplifies a comparison by applying equivalent
transformations as long as possible. The general form is

\begin{plContainer}
\begin{plList}
\Gamma\proofsep&
i.& t_1 \bowtie t_n \leftrightarrow \psi &(\currule)\\
\end{plList}
\end{plContainer}
where $\bowtie$ is as for the proof rule \proofRule{la_generic}, but never
$\simeq$.

The possible transformations are:
\begin{itemize}
    \item $t_1 < t_2 \leftrightarrow \varphi$ where $t_1$ and
    $t_2$ are numerical constants and $\varphi$ is $\top$ if $t_1$ is
    strictly greater than $t_2$ and $\bot$ otherwise.
    \item $t < t \leftrightarrow \bot$
    \item $t_1 \leq t_2 \leftrightarrow \varphi$ where $t_1$ and
    $t_2$ are numerical constants and $\varphi$ is $\top$ if $t_1$ is
    greater than $t_2$ or equal and $\bot$ otherwise.
    \item $t \leq t \leftrightarrow \top$
    \item $t_1 \geq t_2 \leftrightarrow t_2 \leq t_1$
    \item $t_1 < t_2 \leftrightarrow \neg (t_2 \leq t_1)$
    \item $t_1 > t_2 \leftrightarrow \neg (t_1 \leq t_2)$
\end{itemize}
\end{proof-rule}

\begin{proof-rule}{let}{veriT}
  This rule eliminats $\operatorname{let}$. It has the form

  \begin{plContainer}
    \begin{plList}
      \Gamma \proofsep& i_1.& t_{1} \simeq s_{1} &(\dots)\\
      \plLine\\
      \Gamma \proofsep& i_n.& t_{n} \simeq s_{n} &(\dots)\\
    \end{plList}
    \begin{plSubList}
      \plLine\\
      \Gamma, x_1 \mapsto s_1, \dots,  x_n \mapsto s_n,
      \proofsep& j.& u\simeq u' &(\dots)\\
    \end{plSubList}
    \begin{plList}
      \Gamma \proofsep& k.&
      (\operatorname{let} x_1 := t_1, \dots, x_n := t_n. u) \simeq
      u'
      &(\currule;i_1 \dots i_n)\\
    \end{plList}
  \end{plContainer}
  where $\simeq$ is replaced by $\leftrightarrow$ where necessary.

  If for $t_i\simeq s_i$ the $t_i$ and $s_i$ are syntactically equal, the premise
  is skipped.

\end{proof-rule}

\begin{proof-rule}{distinct_elim}{veriT}
This rule eliminates the $\operatorname{distinct}$ predicate. If called with one
argument this predicate always holds:
\begin{plContainer}
\begin{plList}
\proofsep& i.&
(\operatorname{distinct} t) \leftrightarrow \top
&(\currule)\\
\end{plList}
\end{plContainer}

If applied to terms of type $\mathbf{Bool}$ more than two terms can never be
distinct, hence only two cases are possible:

\begin{plContainer}
\begin{plList}
\proofsep& i.&
(\operatorname{distinct} \varphi\;\psi) \leftrightarrow \neg (\varphi \leftrightarrow \psi)
&(\currule)\\
\end{plList}
\end{plContainer}
and
\begin{plContainer}
\begin{plList}
\proofsep& i.&
(\operatorname{distinct} \varphi_1\;\varphi_2\;\varphi_3\dots) \leftrightarrow \bot
&(\currule)\\
\end{plList}
\end{plContainer}

The general case is
\begin{plContainer}
\begin{plList}
\proofsep& i.&
(\operatorname{distinct} t_1 \dots t_n) \leftrightarrow
\bigwedge_{i=1}^{n}\bigwedge_{j=i+1}^{n} t_i\;{\not\simeq}\;t_j
&(\currule)\\
\end{plList}
\end{plContainer}
\end{proof-rule}

\begin{proof-rule}{la_rw_eq}{veriT}
\begin{plContainer}
\begin{plList}
\proofsep& i.&
(t \simeq u) \simeq (t \leq u \land u \leq t)
&(\currule)\\
\end{plList}
\end{plContainer}

\paragraph{Remark.} While the connective could be $\leftrightarrow$,
currently an equality is used.
\end{proof-rule}

\begin{proof-rule}{nary_elim}{veriT}
This rule replaces $n$-ary operators with their equivalent
application of the binary operator. It is never applied to $\land$ or $\lor$.

Thre cases are possible.
If the operator $\circ$ is left associative, then the rule has the form
\begin{plContainer}
\begin{plList}
\Gamma\proofsep&
i.& \bigcirc_{i=1}^{n} t_i \leftrightarrow
(\dots( t_1\circ  t_2) \circ  t_3)\circ \dots  t_n) &(\currule)\\
\end{plList}
\end{plContainer}

If the operator $\circ$ is right associative, then the rule has the form
\begin{plContainer}
\begin{plList}
\Gamma\proofsep&
i.& \bigcirc_{i=1}^{n} t_i \leftrightarrow
( t_1 \circ \dots \circ ( t_{n-2} \circ ( t_{n-1} \circ  t_n)\dots) &(\currule)\\
\end{plList}
\end{plContainer}

If the operator is \emph{chainable}, then it has the form
\begin{plContainer}
\begin{plList}
\Gamma\proofsep&
i.& \bigcirc_{i=1}^{n} t_i \leftrightarrow
( t_1\circ t_2) \land ( t_2 \circ  t_3) \land \dots
\land ( t_{n-1}\circ t_n) &(\currule)\\
\end{plList}
\end{plContainer}
\end{proof-rule}

\begin{proof-rule}{bfun_elim}{veriT}
\begin{plContainer}
\begin{plList}
\proofsep& i.& \psi &(\dots)\\
\proofsep& j.& \varphi &(\currule; i)\\
\end{plList}
\end{plContainer}

The formula $\varphi$ is $\psi$ after boolean functions have been simplified.
This happens in a two step process. Both steps recursively iterate over $\psi$.
The first step expands quantified variable of type $\mathbf{Bool}$. Hence,
$\exists x.t$ becomes $t[\bot/x]\lor t[\top/x]$ and
$\forall x.t$ becomes $t[\bot/x]\land t[\top/x]$. If $n$ variables of sort
$\mathbf{Bool}$ appear in a quantifier, the disjunction (conjunction) has
$2^n$ terms. Each term replaces the variables in $t$ according
to the bits of a number which is increased by one for each subsequent
term starting from zero. The left-most variable corresponds to the
least significant bit.

The second step expands function argument of boolean types by introducing
appropriate if-then-else terms. For example, consider $f(x, P, y)$ where
$P$ is some formula. Then we replace this term by $\operatorname{ite} P\;
f(x, \top, y)\;f(x, \bot, y)$. If the argument is already the constant $\top$
or $\bot$ it is ignored.

\end{proof-rule}

\begin{proof-rule}{ite_intro}{veriT}
Either 
\begin{plContainer}
\begin{plList}
\proofsep& i.& t \simeq (t' \land u_1 \land \dots \land u_n) &(\currule)\\
\end{plList}
\end{plContainer}
or
\begin{plContainer}
\begin{plList}
\proofsep& i.&
\varphi \leftrightarrow (\varphi' \land u_1 \land \dots \land u_n) &(\currule)\\
\end{plList}
\end{plContainer}

The term $t$ (the formula $\varphi$) contains the $\operatorname{ite}$ operator.
Let $s_1, \dots, s_n$ be the terms starting with $\operatorname{ite}$, i.e.
$s_i := \operatorname{ite} \psi_i\;r_i\;r'_i$, then $u_i$ has the form
\begin{equation*}
	\operatorname{ite} \psi_i\;(s_i \simeq r_i)\;(s_i \simeq r'_i)
\end{equation*}
or
\begin{equation*}
	\operatorname{ite} \psi_i\;(s_i \leftrightarrow r_i)\;(s_i \leftrightarrow r'_i)
\end{equation*}
if $s_i$ is of sort $\mathbf{Bool}$. The term $t'$ (the formular
$\varphi'$) is equal to the term $t$ (the formular $\varphi'$) up to the
reordering of equalities where one argument is an $\operatorname{ite}$
term.

\paragraph{Remark.} This rule stems from the introduction of fresh
constants for if-then-else terms inside veriT. Internally $s_i$ is a new
constant symbol and the $\varphi$ on the right side of the equality is
$\varphi$ with the if-then-else terms replaced by the constants. Those
constants are unfolded during proof printing. Hence, the slightly strange
form and the reordering of equalities.
\end{proof-rule}

% \subsection*{Experimental Proof Rules}
% \label{sec:exp-rules}

% \begin{proof-rule}{deep_skolemize}{veriT}
%   This rule is only emitted when the option \texttt{--enable-deep-skolem}
%   is given. This option is experimental and should not be used.
% \end{proof-rule}

%%% Local Variables:
%%% mode: latex
%%% TeX-master: "doc"
%%% End:
